\section{Ethics, trust and decision-making}

\begin{qfdefinition}[Ethics, codes of ethics, professions and trust: why investment management needs standards above the law, and what professionalism requires]\label{def:ethics-trust-professions}
\qffield{Statement} \emph{Ethics} is a set of principles or rules of conduct that guide behaviour when it affects others; widely shared principles are honesty, fairness, diligence, and care and respect for others. Ethical conduct balances self-interest against the direct and indirect consequences for other people \citep[p.~xiii]{standards-practice-handbook-2024}.

\emph{A code of ethics and a profession.} A profession is an occupation whose members hold specialised knowledge that clients cannot easily check, and which therefore binds its members by a published code and enforces it. The CFA Institute Code is aspirational: it states principles such as placing clients' interests and the integrity of the profession above one's own. The Standards of Professional Conduct are the enforceable rules beneath it, applied to members worldwide, regardless of local law or custom (pp.~v, xviii). A code becomes a ``dusty compilation'' unless it is built into the firm's culture, training and incentives (p.~xvi).

\emph{Trust and the need for high standards.} Capital markets channel savings to productive uses only if investors trust that markets are fair and that professionals act for them. Investors who distrust them invest less or demand a higher risk premium, which raises the cost of capital (p.~xiv). Gennaioli, Shleifer and Vishny model the other side: investors who trust a manager perceive delegated investments as less risky, take more risk than they would alone and accept fees that do not fall to cost. Trust is valuable but exploitable. Managers who are trusted can charge more and gain from pandering to clients' mistaken beliefs rather than correcting them \citep[NBER WP 18174, pp.~1--5]{gennaioli-shleifer-vishny-money-doctors-2012}. Information asymmetry and delegated discretion are why investment work needs standards stricter than those of an ordinary sale.

\emph{Professionalism} means acting with competence and integrity in the client's interest: maintaining the knowledge one's role requires (Standard~I(E)), exercising independent judgement, communicating plainly and placing client interests first (Code, p.~xviii; Standard~III(A), p.~65).

\qffield{Practice} \emph{Why the standard matters} (original illustration). An adviser recommends a fund that pays her firm a 0.50 percent trail. If clients cannot verify the advice, disclosure and a duty of loyalty are the only protections; once a mis-selling case becomes public, every client discounts every adviser's advice. In Gennaioli, Shleifer and Vishny's model, trust raises fees even without better performance, so the profession's code is what keeps trust earned rather than merely exploited (Interpretation).

\qffield{Conditions} The Code and Standards bind CFA Institute members and candidates; firms may adopt them voluntarily (Handbook, p.~x). The model of trust is a theoretical account with supporting evidence cited by its authors, not a measurement of any market.

\qffield{Pitfalls} Treating legal compliance as the whole of ethics. Treating a firm's written code as evidence of an ethical culture. Assuming that clients' trust shows the advice is good.

\qffield{See also} \hyperref[meth:ethical-decision-framework]{decision framework}; \hyperref[def:code-of-ethics-standards]{the Code and Standards}.

\end{qfdefinition}

\begin{qfmethod}[Ethical challenges and decision-making: situational influences, ethical against legal standards, and the identify--consider--act--reflect framework applied to a case]\label{meth:ethical-decision-framework}
\qffield{Statement} \emph{Challenges to ethical behaviour.} Many decisions are not clearly right or wrong, and even well-intentioned professionals can be tempted to cut corners: situational influences can overpower the best of intentions \citep[p.~xvii]{standards-practice-handbook-2024}. Pressure to conform to a group, obedience to authority, incentives tied to short-term results and overconfidence in one's own judgement all distort decisions \citep[p.~1]{ethical-decision-framework-2014}. Compensation that rewards revenue regardless of how it is earned is the standard example (Handbook, pp.~xv, 124--125).

\emph{Ethical against legal standards.} Laws and ethics both constrain self-interest that could harm others, but they differ. Laws state what is required, and set a minimum that is often enforced with limited resources. Ethics asks what is right, and can require more. Conduct may be legal but unethical, such as a conflict disclosed in boilerplate the client cannot understand. Conversely, protecting clients from an employer's illegal conduct can justify acts that would otherwise breach the duty to the employer, such as contradicting instructions or keeping copies of records (Handbook, pp.~xv, 110). Where law and the Standards differ, members follow the stricter (Standard~I(A), p.~2). US advisers' own legal duty shows the gap is narrowing. Under the SEC's 2019 interpretation, disclosure and consent allow a conflict but do not satisfy the duty to act in the client's best interest, and saying an adviser ``may'' have a conflict that it actually has is inadequate \citep[84 FR 33669, at 33672, 33676]{sec-adviser-standard-of-conduct-2019}.

\emph{The framework} \citep[p.~1]{ethical-decision-framework-2014} is a structure for situations without a clear right and wrong, not a checklist:
\begin{itemize}
\item \emph{Identify} the principles at issue (fair dealing, disclosure, loyalty, diligence), the duties owed and to whom, the facts known and still needed, and any conflicts of interest.
\item \emph{Consider} situational influences and internal biases, alternative actions, and outside guidance (compliance, legal, a trusted colleague).
\item \emph{Act}: decide, or elevate the issue to a higher authority and follow up.
\item \emph{Reflect} on what was learned and on one's strengths and weaknesses.
\end{itemize}

\qffield{Practice} \emph{Corin Asset Management} (original). A portfolio manager, Ines Vale, is told by her head of sales to keep a large client's underperforming account out of the firm's equity composite until it is ``fixed'', because the composite is in a pitch to a pension fund next week.
\begin{itemize}
\item \emph{Identify}: duty to prospective clients not to misrepresent performance (Standards~I(C), III(D)); duty to the employer; the facts needed are whether the account meets the composite's definition. Conflict: her bonus depends on new mandates.
\item \emph{Consider}: obedience to a senior manager and the group's eagerness to win are pressures. Alternatives: comply; refuse silently; raise with compliance; present the composite with the account and explain the underperformance.
\item \emph{Act}: the account meets the definition, so excluding it would misstate the record. She explains this to the head of sales and, when he insists, takes the issue to the compliance officer, who rules that the composite must include it. She is not required to resign; she must not distribute a record she knows is inflated (Standard~I(A) dissociation, pp.~3, 6).
\item \emph{Reflect}: the firm lacked a written composite-definition policy; she proposes one.
\end{itemize}
The legal question (would a regulator act?) and the ethical one (is the presentation fair?) give the same answer here only because the stricter standard is applied.

\qffield{Conditions} The framework is CFA Institute's (2014 card). The SEC interpretation applies to US investment advisers and is cited only to contrast a legal fiduciary standard with the Code.

\qffield{Pitfalls} Treating ``not illegal'' as ``acceptable''. Escalating without following up. Relying on a lawyer's view that a regulator could not prove wrongdoing as a reason to proceed (Handbook, p.~6).

\qffield{See also} \hyperref[def:ethics-trust-professions]{ethics and trust}; \hyperref[meth:evaluating-conduct-cases]{evaluating conduct}.

\end{qfmethod}
